Real-time analysis of public instrument performance videos requires models that remain useful when audio and visual evidence are weakly synchronized, partially occluded, or only coarsely labeled. Existing multimodal music research has made strong progress in retrieval, correspondence learning, question answering, and large-model music understanding, yet it has only partially addressed clip-level instrument classification under tight latency budgets. This paper presents \method{} (Audio-visual Instrument Motion-Enhanced Recognizer), a lightweight audio-visual framework for real-time instrument classification on public performance videos. The model combines a compact log-mel audio encoder, a motion-aware visual encoder, and a latency-aware fusion gate that can down-weight unreliable modalities when performer motion and soundtrack evidence disagree. In addition, we use a coarse phase label only as auxiliary supervision and as a diagnostic variable rather than as a standalone benchmark target. The evaluation protocol reports clean-set recognition, batch-one inference latency, audio-visual offset robustness, and modality corruption on URMP and Solos. On the evaluated setup, \method{} reaches Macro-F1 scores of 0.84 on URMP and 0.76 on Solos, improving over naive late fusion by 0.03 and 0.04 while reducing latency from 25 to 18 ms per clip and lowering the offset-induced drop from 0.07 to 0.03. These results support \method{} as a lightweight instrument-recognition model with phase-aware diagnostics on public performance videos.
